% from ega4-2-fr.tex
\oldpage[IV-2]{141}

idéal premier minimal de $A$. Posons pour simplifier
$A'=\widehat A$, $M'=\widehat M=M\otimes_AA'$ $(0_{\mathrm I}, 7.3.3)$; on
sait (0, 16.5.2) que $M'$ est un $A'$-module de Cohen-Macaulay, ce qui entraîne
que pour tout idéal premier $\mathfrak p'$ de $A'$,
$M'_{\mathfrak p'}$ est un $A'_{\mathfrak p'}$-module de Cohen-Macaulay
(0, 16.5.10). Compte tenu de (I, 3.6.5), on voit que si l'on pose
$A''=A'\otimes_AA_{\mathfrak p}$ et
$M''=M'\otimes_AA_{\mathfrak p}$, $M''$ est un $A''$-module de
Cohen-Macaulay. Pour tout idéal premier $\mathfrak p''$ de $A''$ au-dessus
de $\mathfrak p$,
$M''_{\mathfrak p''}=M_{\mathfrak p}\otimes_{A_{\mathfrak p}}
A''_{\mathfrak p''}$ est donc un $A''_{\mathfrak p''}$-module de
Cohen-Macaulay (0, 16.5.10); d'autre part, $M_{\mathfrak p}$ est par hypothèse
un $A_{\mathfrak p}$-module de Cohen-Macaulay et
$A''_{\mathfrak p''}$ est un $A_{\mathfrak p}$-module plat puisque $A'$ est
un $A$-module plat $(0_{\mathrm I}, 7.3.3\text{ et }6.3.2)$. On conclut de (6.3.3)
que, si $k$ est le corps résiduel de $A_{\mathfrak p}$,
$k\otimes_{A_{\mathfrak p}}A''_{\mathfrak p''}$ est un anneau de
Cohen-Macaulay. Mais puisque $A_{\mathfrak p}$ est de dimension $0$, les
idéaux premiers de $A''$ correspondent biunivoquement à ceux de
$k\otimes_{A_{\mathfrak p}}A''$ (I, 3.5.7), et si $\mathfrak q''$ est l'idéal
premier de $k\otimes_{A_{\mathfrak p}}A''$ correspondant à
$\mathfrak p''$, les anneaux locaux
$(k\otimes_{A_{\mathfrak p}}A'')_{\mathfrak q''}$ et
$k\otimes_{A_{\mathfrak p}}A''_{\mathfrak p''}$ sont isomorphes. Par suite
(0, 16.5.13), l'anneau $k\otimes_{A_{\mathfrak p}}A''$ est un anneau de
Cohen-Macaulay. C.Q.F.D.

\subsection{\texorpdfstring{Platitude et propriété $(S_k)$}{Platitude et propriété (S\string_k)}}

\begin{proposition}[6.4.1]
\label{IV.2.s6.4.1-fr}
Soient $X$, $Y$ deux préschémas localement noethériens, $f:X\to Y$ un
morphisme, $\mathscr E$ un $\mathcal O_Y$-Module cohérent,
$\mathscr F$ un $\mathcal O_X$-Module cohérent et $f$-plat.
\begin{enumerate}
\item[(i)] Soit $x\in\operatorname{Supp}(\mathscr F)$ tel que
$\mathscr E\otimes_Y\mathscr F$ ait la propriété $(S_k)$ au point $x$; alors
$\mathscr E$ a la propriété $(S_k)$ au point $y=f(x)$.
\item[(ii)] Supposons que pour tout
$y\in f(\operatorname{Supp}(\mathscr F))$,
$\mathscr F_y=\mathscr F\otimes_{\mathcal O_y}k(y)$ ait la propriété
$(S_k)$; alors, si pour un point
$y\in f(\operatorname{Supp}(\mathscr F))$, $\mathscr E$ a la propriété
$(S_k)$ au point $y$, $\mathscr E\otimes_Y\mathscr F$ a la propriété
$(S_k)$ en tout point de $f^{-1}(y)$.
\end{enumerate}
\end{proposition}

(i) On sait $(\mathbf{Err}_{\mathrm{III}},30)$ qu'il existe un sous-préschéma
fermé $X'$ de $X$ ayant pour espace sous-jacent
$\operatorname{Supp}(\mathscr F)$ et un $\mathcal O_{X'}$-Module cohérent
$\mathscr F'$ tels que $\mathscr F=j_*(\mathscr F')$, où $j$ est l'injection
canonique $X'\to X$; on peut remplacer $X$ par $X'$ et $\mathscr F$ par
$\mathscr F'$, autrement dit supposer que
$\operatorname{Supp}(\mathscr F)=X$. Soit alors $y'$ une générisation de $y$
dans $Y$; on sait (2.3.4) qu'il y a une générisation $x'$ de $x$ dans $X$
telle que $y'=f(x')$; on peut en outre supposer que $x'$ est point générique
d'une composante irréductible de $f^{-1}(y')$; on a par suite
$\dim(\mathscr F_{x'}\otimes_{\mathcal O_{y'}}k(y'))=0$. En vertu de (6.1.2),
on a donc, en posant $\mathscr G=\mathscr E\otimes_Y\mathscr F$,
\[
\dim(\mathscr G_{x'})=\dim(\mathscr E_{y'})
\]
et en vertu de (6.3.1)
\[
\operatorname{prof}(\mathscr G_{x'})=\operatorname{prof}(\mathscr E_{y'})
\]
compte tenu de ce que la profondeur est toujours au plus égale à la
dimension. Par hypothèse, on a
\[
\operatorname{prof}(\mathscr G_{x'})
\geq\inf(k,\dim(\mathscr G_{x'}))
\]
donc
\[
\operatorname{prof}(\mathscr E_{y'})
\geq\inf(k,\dim(\mathscr E_{y'}))
\]
ce qui démontre la première assertion.

